\section{Orientations}

\subsection{Orientations of manifolds}

\bdefn

    Two bases $(E_1,\dots,E_n)$ and $(\tilde E_1,\dots,\tilde E_n)$ of a vector space $V$ are {\emphcolor consistently oriented} if their transformation matrix has
    positive determinant.
    That is, the matrix $(B^j_i)$ given by $E_i=B^j_i\tilde E_j$ has positive determinant.
    This is clearly an equivalence relation with two equivalence classes.

    An {\emphcolor orientation} for $V$ is a choice of one of these equivalence classes.
    For a basis $(E_1,\dots,E_n)$, denote by $[E_1,\dots,E_n]$ its equivalence class and $-[E_1,\dots,E_n]=[-E_1,\dots,E_n]$ the other equivalence class.
    Any basis in this orientation is said to be {\emphcolor positively oriented}, or just {\emphcolor oriented}, and all others are {\emphcolor negatively oriented}.

    An orientation of the zero-dimensional space is a choice of one of $\pm1$.

\edefn

\bexam

    The {\emphcolor standard orientiation} of $\bR^n$ is $[e_1,\dots,e_n]$.

\eexam

\bprop

    Let $V$ be an $n$-dimensional vector space.
    Then every nonzero $\omega\in\Lambda^n(V^*)$ determines an orientation $\O_\omega$ of $V$, where $(E_1,\dots,E_n)$ is oriented iff $\omega(E_1,\dots,E_n)>0$.
    Two nonzero $n$-covectors determine the same orientation iff they are positive multiples of one another.

\eprop

\bproof

    The zero-dimensional case is immediate since $\Lambda^0(V^*)\cong\bR$.
    Otherwise let $\omega\in\Lambda^n(V^*)$, and $(E_i)$ and $(\tilde E_i)$ be bases.
    We claim that they are consistently oriented iff $\omega(E_1,\dots,E_n)$ and $\omega(\tilde E_1,\dots,\tilde E_n)$ have the same sign.
    Let $B=(B^j_i)$ be their transition matrix so
    $$ \omega(\tilde E_1,\dots,\tilde E_n) = \omega(BE_1,\dots,BE_n) = \det(B)\omega(E_1,\dots,E_n) . $$
    So they have the same sign iff $\det(B)>0$, meaning the bases are consistently ordered as desired.
    The last claim is clear.
    \qed

\eproof

We say that a covector $\omega\in\Lambda^n(V^*)$ which determines the orientation of $V$ is an {\emphcolor (positively) oriented covector}.
For example, $e_1\wedge\cdots\wedge e_n$ is a positively oriented covector for the standard orientation of $\bR^n$.

Now we can discuss orientations of manifolds.

\bdefn

    A {\emphcolor (pointwise) orientation} of a smooth manifold is a choice of orientation of $T_pM$ for each $p\in M$.
    If $M$ is endowed with a pointwise orientation, a local frame $(E_i)$ is {\emphcolor (positively) oriented} if $(E_i|_p)$
    is a positively oriented basis for each $p$ in the frame's domain.
    We similarly define a negatively oriented frame.

    A pointwise orientation is {\emphcolor continuous} if every point of $M$ is in the domain of an oriented local frame.
    An {\emphcolor orientation} of $M$ is a continuous pointwise orientation.
    A manifold is said to be {\emphcolor orientable} if it possesses a continuous pointwise orientation, otherwise it is
    {\emphcolor nonorientable}.

    An {\emphcolor oriented manifold} is a pair $(M,\O)$ where $\O$ is an orientation on $M$.
    Denote the orientation of $T_pM$ by $\O_p$.

\edefn

This definition is valid for manifolds with boundary.
For a zero-dimensional manifold, a pointwise orientation is a choice of $\pm1$ for every point $p\in M$.
In the zero-dimensional case every pointwise orientation is continuous.

Note that if $M$ is oriented, any local frame defined on a connected domain is either positively or negatively oriented.
Let $(E_i)$ be a local frame on a connected domain $U$, then define $U\to\bR$ by mapping $p\in U$ to $+1$ if $(E_i|_p)$ is
positively oriented and $-1$ otherwise.
Then this is locally constant: for $p\in U$ take a positively oriented frame $(F_i)$ defined in a neighborhood of $p$,
which we can restrict to a connected open subset of $U$.
Let $B$ be the transition map between $E$ and $F$, then $E_p$ is positively oriented iff $\det(B_p)>0$, so in this open
subset the map is equal to $\sgn\det(B)$.
Since $\det(B)$ is nonzero, continuous, and the domain is connected, we have that its sign is constant.
So the map is locally constant, and since the domain is connected it must be constant.

\bprop

    Let $M$ be a smooth $n$-manifold, then any nonvanishing $n$-form $\omega$ on $M$ determines a unique orientation for
    which $\omega$ is positively oriented at each point.
    Conversely if $M$ is given an orientation, there is a smooth nonvanishing $n$-form on $M$ that is positively oriented
    at each point.

\eprop

\bproof

    Let $\omega\in\Omega^n(M)$ be nonvanishing.
    Then $\omega_p$ defines an orientation on $T_pM$ for all $p\in M$, we check that this orientation is continuous.
    This is trivial for $n=0$, so we assume $n\geq1$.

    By the above discussion, every local frame on a connected neighborhood should be positively or negatively oriented.
    So let $p\in M$ and take a local frame on a connected neighborhood $U$ of $p$, call it $(E_i)$ with dual coframe
    $(\epsilon^j)$.
    Write $\omega=f\epsilon^1\wedge\cdots\wedge\epsilon^n$, and since $\omega$ is nonvanishing $f$ is nonzero over all of $U$
    so
    $$ \omega(E_1,\dots,E_n) = f \neq 0 . $$
    Since $U$ is connected, $f$ is either always positive or always negative.
    If $f$ is always positive then $(E_i)$ is positively oriented and we have finished.
    Otherwise $f$ is always negative so $\omega(-E_1,E_2,\dots,E_n)=-f$ is always positive and we take the local frame
    $(-E_1,E_2,\dots,E_n)$ and we have finished.

    For $p\in M$ we take a connected neighborhood $U_p$ with positively oriented frame $(E_i^p)$.
    Let $(\epsilon^j_p)$ be its dual coframe and define the local $n$-form $\omega^p=\epsilon^1_p\wedge\cdots\epsilon^n_p$.
    Then let $\psi_p$ be a partition of unity subordinate to $\set{U_p}_p$ and define
    $$ \omega = \sum_p\psi_p\omega^p . $$
    This is nonvanishing since for any $(E_i^p)$, $\omega^q(E_i^p)\geq0$ for all $q$ since it is positively oriented.
    Since $\psi_p(p)=1$ this implies $\omega(E_i^p)>0$.
    And since $\omega$ is a positive linear combination of positively oriented $n$-forms at each point it is positively
    oriented at each point.
    \qed

\eproof

\bdefn

    A nowhere-vanishing $n$-form on $M$ is called an {\emphcolor orientation form}.
    If $M$ is oriented, an $n$-form giving its orientation is called {\emphcolor (positively) oriented}.

\edefn

Two orientation forms $\omega,\tilde\omega$ are easily seen to give the same orientation iff $\omega=f\tilde\omega$ for some
positive smooth real-valued $f$.

\bdefn

    A smooth coordinate chart is {\emphcolor (positively) oriented} if its coordinate frame $(\ipdv{x^i})$ is positively
    oriented, and {\emphcolor negatively oriented} if it the coordinate frame is negatively oriented.

    An atlas $\set{(U_i,\phi_i)}_i$ is {\emphcolor consistently oriented} iff the transition maps $\phi_j\circ\phi_i^{-1}$
    have positive determinant for all $i,j$.

\edefn

\bprop

    Let $M$ be a smooth positive-dimensional manifold.
    Any consistently oriented smooth atlas for $M$ uniquely determines an orientation for which the charts are
    all positively oriented.

    Conversely if $M$ is oriented, and either $\partial M=\emptyset$ or $\dim M>1$, then the collection of all oriented smooth
    charts is a consistently oriented atlas for $M$.

\eprop

\bproof

    Let $(U_i,\phi_i)$ be a consistently oriented smooth atlas.
    Then each chart determines a pointwise orientation of its domain, and consistency proves that these pointwise orientations
    agree on their intersections.
    Since the chart covers the manifold, this is a continuous orientation.

    If $M$ is oriented and $\partial M=\emptyset$ or $\dim M>1$, each point in $M$ is in the domain of a connected smooth chart.
    If the domain is negatively oriented, we swap $x^1$ with $-x^1$ to obtain a positively oriented chart.
    As they are all positively oriented, their transition functions must have positive Jacobian.

    For a boundary chart, it maps into $\bH^n$ so the last coordinate must be positive.
    For $\dim M=1$ this is the only coordinate so we cannot swap its sign.
    \qed

\eproof

The following proposition is easily seen.

\bprop

    If $M_1,\dots,M_k$ are orientable smooth manifolds, then there is a unique orientation on $M=\prod_iM_i$
    called the {\emphcolor product orientation}, such that if $\omega_i$ is an orientation form for $M_i$
    then $\pi_1^*\omega_1\wedge\cdots\wedge\pi_n^*\omega_k$ is an orientation form for the product orientation.

\eprop

\bdefn

    Let $M,N$ be oriented smooth manifolds with positive dimension.
    A local diffeomorphism $F\colon M\to N$ is {\emphcolor orientation preserving} if for each $p\in M$ its differential
    $dF_p$ maps oriented bases of $T_pM$ to oriented bases of $T_{Fp}N$.
    Similarly it is {\emphcolor orientation reversing} if it maps negatively oriented bases to oriented bases.

    For zero-dimensional manifolds, $F$ is orientation preserving if for every $p\in M$, both $p$ and $F(p)$ have the same
    orientation (choice of $\pm1$).

\edefn

